
\section{ผลเปรียบเทียบ}

ตารางเปรียบเทียบค่าความสูญเสียและตัวชี้วัดระหว่างการฝึกแบบจำลองของ YOLOv26n และโมเดลตัวอื่นๆ


\begin{table}[H]
\centering
\small
\caption{เปรียบเทียบประสิทธิภาพของโมเดลตรวจจับใบหน้า YOLOv26n, YuNet, RetinaFace และ MTCNN บนชุดข้อมูล WIDER FACE โดยใช้ตัวชี้วัด AP (Average Precision) และ Latency (เวลาในการประมวลผลต่อภาพ) ที่ความละเอียด 320x320 และ 640x640 ที่ทำงานบน ONNX}
\label{tab:yunet_comparison}
\resizebox{\textwidth}{!}{%
\begin{tabular}{c l r r c c c c}
\toprule
\textbf{Input size} &
\textbf{Methods} &
\textbf{\#FLOPs (M)} &
\textbf{$AP_{easy}$} &
\textbf{$AP_{medium}$} &
\textbf{$AP_{hard}$} &
\textbf{Latency (ms)}
\\
\midrule

% ==================== 320 x 320 ====================
\multirow{4}{*}{320 $\times$ 320}

& YOLOv26n (Ours)
& 194
& 0.894
& 0.8677
& 0.7257
& 7.2
\\

& YuNet
& 149
& 0.836
& 0.747
& 0.395
& 2.2
\\

& RetinaFace
& 11\,070
& 0.868
& 0.742
& 0.341
& 49.1
\\

& MTCNN
& 3\,359
& 0.923
& 0.862
& 0.504
& 17.3
\\

\midrule

% ==================== 640 x 640 ====================
\multirow{4}{*}{640 $\times$ 640}

& YOLOv26n (Ours)
& 692
& 0.907
& 0.894
& 0.880
& 17.8
\\

& YuNet
& 595
& 0.899
& 0.869
& 0.691
& 20.4
\\

& RetinaFace
& 44\,260
& 0.943
& 0.908
& 0.659
& 232.7
\\

& MTCNN
& 11\,545
& 0.949
& 0.935
& 0.814
& 104.0
\\

\midrule
\bottomrule
\end{tabular}%
}
\end{table}

จากผลลัพธ์ตามตาราง \ref{tab:yunet_comparison} จะเห็นว่า YOLOv26n มีค่า AP โดยรวมที่สูงและมี Latency ที่ต่ำเมื่อเทียบกับโมเดลตัวอื่น ทำให้ YOLOv26n เป็นตัวเลือกที่ดีสำหรับการตรวจจับใบหน้าในสภาพแวดล้อมที่ต้องการความเร็วและความแม่นยำ มากกว่า YuNet ที่เหมาะสำหรับงานภาพหรือวิดีโอที่มีขนาดเล็ก



\section{สรุปผลการดำเนินโครงงาน}

โครงงานนี้พัฒนาระบบ Blur System for Privacy Protection เพื่อช่วยให้ผู้ใช้ปกปิดข้อมูลระบุตัวบุคคลในภาพและวิดีโอก่อนนำสื่อไปใช้งานต่อ จากการตรวจสอบซอร์สโค้ดในรอบปัจจุบัน ระบบมีกระบวนการทำงานหลักตั้งแต่การอัปโหลดสื่อ การตรวจจับและจัดกลุ่มใบหน้า การเลือกบุคคลจาก Face Map การกำหนดรูปแบบการปกปิด การวาดและติดตามกรอบด้วยตนเอง การดูตัวอย่าง และการส่งออกไฟล์

\begin{table}[H]
\centering
\caption{สรุปผลเทียบกับวัตถุประสงค์}
\label{tab:objective-summary}
\begin{tabular}{|p{0.10\textwidth}|p{0.43\textwidth}|p{0.31\textwidth}|}
\hline
\textbf{ข้อ} & \textbf{วัตถุประสงค์} & \textbf{ผลและสถานะ} \\
\hline
1 & พัฒนาเว็บแอปพลิเคชันสำหรับตรวจจับและปกปิดใบหน้า & มีส่วน API ของ frontend และ backend และการสร้างผลลัพธ์ บรรลุในระดับต้นแบบ \\
\hline
2 & จัดกลุ่มใบหน้าและเลือกปกปิดรายบุคคล & มีเวกเตอร์ลักษณะจาก InsightFace/ArcFace, Face Map และสถานะการเลือกปกปิด บรรลุในระดับต้นแบบ \\
\hline
3 & รองรับการปกปิดด้วยตนเองและติดตามวัตถุ & มี Manual Control และ KCF/CSRT เป็นทางเลือกสำรอง บรรลุในระดับต้นแบบ \\
\hline
4 & เลือกเอฟเฟกต์ ดูตัวอย่าง และส่งออกผลลัพธ์ & รองรับ Gaussian Blur, Pixelation, Black Box การดูตัวอย่าง และการส่งออก บรรลุในระดับต้นแบบ \\
\hline
5 & ประเมินความแม่นยำ ความเร็ว และความครอบคลุม & มีแผนและตัวชี้วัด แต่ยังไม่มีผลการประเมินที่ตรวจสอบได้ อยู่ระหว่างดำเนินการ \\
\hline
\end{tabular}
\end{table}

\section{ข้อจำกัดของระบบ}

\subsection{ข้อจำกัดด้านการทำงาน}

ระบบตรวจจับอัตโนมัติเน้นใบหน้าเป็นหลัก จึงยังไม่ครอบคลุมป้ายทะเบียน ข้อความ ลายเซ็น หรือข้อมูลส่วนบุคคลประเภทอื่นโดยอัตโนมัติ ผู้ใช้ต้องวาดกรอบ Manual Blur เพื่อปกปิดข้อมูลเพิ่มเติม ความต่อเนื่องของการจัดกลุ่มบุคคลและการติดตามกรอบอาจลดลงเมื่อภาพเบลอมาก วัตถุเคลื่อนที่เร็ว ถูกบดบัง หรือมีมุมหน้าเปลี่ยนแปลงมาก

สถานะของงานเบื้องหลังถูกเก็บไว้ในหน่วยความจำ จึงไม่เหมาะกับการทำงานหลาย instance หรือการรีสตาร์ตเซิร์ฟเวอร์ นอกจากนี้ backend จะกำหนดตัวแปรขนาดไฟล์สูงสุดไว้ แต่ยังไม่ได้ใช้บังคับก่อนบันทึกไฟล์ และในปัจจุบันยังไม่มีกระบวนการลบไฟล์อัปโหลดตามเวลา

\section{ปัญหา อุปสรรค และแนวทางแก้ไข}

\begin{table}[H]
\centering
\caption{ปัญหา อุปสรรค และแนวทางแก้ไข}
\label{tab:problems}
\begin{tabular}{|p{0.25\textwidth}|p{0.29\textwidth}|p{0.30\textwidth}|}
\hline
\textbf{ปัญหาที่พบ} & \textbf{ผลกระทบ} & \textbf{แนวทางแก้ไข} \\
\hline
ภาพคุณภาพต่ำ แสงน้อย หรือใบหน้าถูกบดบัง & การตรวจจับและจัดกลุ่มบุคคลไม่เสถียร & เก็บชุดทดสอบหลายสภาพ ปรับค่าเกณฑ์ และเพิ่มเครื่องมือแก้ไขด้วยตนเอง \\
\hline
วัตถุส่วนตัวที่ไม่ใช่ใบหน้า & ไม่ถูกปกปิดอัตโนมัติ & ใช้ Manual Control ในปัจจุบัน และพัฒนาโมเดลป้ายทะเบียน/OCR ในระยะต่อไป \\
\hline
วิดีโอความละเอียดสูงหรือมีความยาวมาก & เวลาในการสร้างผลลัพธ์และการใช้หน่วยความจำสูง & เพิ่มคิวงาน สร้างตัวอย่างความละเอียดต่ำ และวัดประสิทธิภาพตามความละเอียด \\
\hline
Type safety ของส่วนหน้ายังไม่สมบูรณ์ & การตรวจ lint ไม่ผ่านและเสี่ยงต่อความผิดพลาดขณะเชื่อม API & แทนที่ \texttt{any} ด้วย type/interface และแก้ผลตรวจ lint เป็นลำดับแรก \\
\hline
สถานะงานเก็บในหน่วยความจำ & สถานะอาจสูญหายเมื่อส่วนหลังรีสตาร์ต & ย้ายสถานะและคิวงานไปยังฐานข้อมูลหรือ message queue \\
\hline
\end{tabular}
\end{table}

\section{ข้อเสนอแนะในการพัฒนาต่อไป}

\subsection{การปรับปรุงฟังก์ชันการทำงาน}

ควรเพิ่มแบบจำลองตรวจจับป้ายทะเบียนและโมดูล OCR/PII เพื่อลดการวาดกรอบด้วยตนเอง เพิ่มหน้าสำหรับแก้ไขผลการจัดกลุ่มบุคคล และอนุญาตให้กำหนดกรอบเริ่มต้นได้มากกว่าหนึ่งจุดในวิดีโอ นอกจากนี้ควรเพิ่มระบบลบไฟล์ชั่วคราว กำหนดขนาดไฟล์และสิทธิ์เข้าถึงให้ชัดเจน และปรับปรุงคุณภาพการตรวจจับและติดตามวัตถุหรือใบหน้าทั้งในด้านความเร็วและความแม่นยำเพิ่มเติม

\subsection{แนวทางการวิจัยต่อยอด}

จัดเตรียมชุดทดสอบที่มีสภาพแสง มุมหน้า การเคลื่อนไหว ความละเอียด และจำนวนบุคคลต่างกัน แล้ววัด precision, recall, identity consistency, blur coverage และเวลาเฉลี่ยต่อเฟรม นอกจากนี้ทดสอบว่าระดับการเบลอแต่ละแบบเพียงพอต่อการลดความสามารถในการระบุตัวบุคคลหรือไม่ และประเมินประสบการณ์ผู้ใช้ด้วยแบบสอบถามหลังใช้งาน
